\chapter{Chapter05 — FKN Optimality}
%%% generated by the blueprint pipeline from FABL.Chapter05.FKNOptimality
\section{Overview}

Declarations of the Lean module \texttt{FABL.Chapter05.FKNOptimality}.

\section{Declarations}

\begin{definition}[Gaussian mass parameter]
\lean{FABL.fknOptimalityGaussianMass}
\label{FABL.fknOptimalityGaussianMass}
\leanok
For each natural number $q$, define the Gaussian mass parameter to be $1/(q+2)$, which
lies strictly between $0$ and $1$.
\end{definition}

\begin{definition}[Gaussian threshold]
\lean{FABL.fknOptimalityThreshold}
\label{FABL.fknOptimalityThreshold}
\leanok
\uses{FABL.fknOptimalityGaussianMass, FABL.standardGaussianUpperQuantile}
For each natural number $q$, define the threshold as the standard Gaussian upper-tail
quantile corresponding to mass $1/(q+2)$.
\end{definition}

\begin{definition}[Gaussian boundary value]
\lean{FABL.fknOptimalityGaussianBoundary}
\label{FABL.fknOptimalityGaussianBoundary}
\leanok
\uses{FABL.fknOptimalityGaussianMass, FABL.gaussianIsoperimetric}
For each natural number $q$, define the Gaussian boundary value at mass $1/(q+2)$;
equivalently, it is the standard Gaussian density evaluated at the standard Gaussian
upper-tail quantile for that mass.
\end{definition}

\begin{theorem}[Tail at the optimality threshold]
\lean{FABL.standardGaussianUpperTail_fknOptimalityThreshold}
\label{FABL.standardGaussianUpperTail_fknOptimalityThreshold}
\leanok
\uses{FABL.fknOptimalityGaussianMass, FABL.fknOptimalityThreshold, FABL.standardGaussianUpperTail, FABL.standardGaussianUpperTail_quantile}
\difficulty{1}
For every natural number $q$, the standard Gaussian upper-tail probability at the
threshold corresponding to Gaussian mass $1/(q+2)$ is exactly $1/(q+2)$.
\end{theorem}

\begin{theorem}[Gaussian boundary as a density]
\lean{FABL.fknOptimalityGaussianBoundary_eq_density}
\label{FABL.fknOptimalityGaussianBoundary_eq_density}
\leanok
\uses{FABL.fknOptimalityGaussianBoundary, FABL.fknOptimalityGaussianMass, FABL.fknOptimalityThreshold, FABL.gaussianIsoperimetric_apply_of_mem_Ioo}
\difficulty{1}
For every $q\in\mathbb{N}$, the Gaussian boundary associated with $q$ equals the
standard Gaussian density evaluated at the standard Gaussian upper quantile
corresponding to mass $1/(q+2)$.
\end{theorem}

\begin{theorem}[Positive Gaussian boundary]
\lean{FABL.fknOptimalityGaussianBoundary_pos}
\label{FABL.fknOptimalityGaussianBoundary_pos}
\leanok
\uses{FABL.fknOptimalityGaussianBoundary, FABL.fknOptimalityGaussianBoundary_eq_density, FABL.fknOptimalityThreshold}
\difficulty{3}
For every $q\in\mathbb{N}$, the Gaussian boundary associated with $q$ is strictly
positive.
\end{theorem}

\begin{theorem}[Vanishing Gaussian mass parameter]
\lean{FABL.tendsto_fknOptimalityGaussianMass_zero}
\label{FABL.tendsto_fknOptimalityGaussianMass_zero}
\leanok
\uses{FABL.fknOptimalityGaussianMass}
\difficulty{3}
For each natural number $q$, set $m_q=1/(q+2)$. Then $m_q$ tends to $0$ as $q$ tends to
infinity.
\end{theorem}

\begin{theorem}[Convergence of the Gaussian mass to zero]
\lean{FABL.tendsto_fknOptimalityGaussianMass_atBot}
\label{FABL.tendsto_fknOptimalityGaussianMass_atBot}
\leanok
\uses{FABL.fknOptimalityGaussianMass, FABL.tendsto_fknOptimalityGaussianMass_zero}
\difficulty{3}
As the natural number $q$ tends to infinity, the quantity $1/(q+2)$ tends to zero within
the interval $(0,1)$, equipped with its order topology.
\end{theorem}

\begin{definition}[Approximation by a Hamming tail]
\lean{FABL.IsFKNOptimalityApproximation}
\label{FABL.IsFKNOptimalityApproximation}
\leanok
\uses{FABL.fknOptimalityGaussianBoundary, FABL.fknOptimalityGaussianMass, FABL.fknOptimalityThreshold, FABL.fourierWeightAtLevel, FABL.hammingUpperTailIndicator}
For natural numbers $q$ and $r$, the approximation property holds when both of the
following inequalities hold on the uniformly distributed sign cube of dimension $r+1$:
the expected value of the Hamming upper-tail indicator at the threshold for $q$ differs
from $1/(q+2)$ by less than half that mass, and its level-one Fourier weight differs
from the square of the Gaussian boundary value for $q$ by less than half that square.
\end{definition}

\begin{theorem}[Approximation to Gaussian boundary data]
\lean{FABL.exists_fknOptimalityApproximation}
\label{FABL.exists_fknOptimalityApproximation}
\leanok
\uses{FABL.IsFKNOptimalityApproximation, FABL.fknOptimalityGaussianBoundary, FABL.fknOptimalityGaussianBoundary_eq_density, FABL.fknOptimalityGaussianBoundary_pos, FABL.fknOptimalityGaussianMass, FABL.fknOptimalityThreshold, FABL.fourierWeightAtLevel, FABL.hammingUpperTailIndicator, FABL.proposition5_25, FABL.standardGaussianUpperTail_fknOptimalityThreshold}
\difficulty{4}
For each natural number $q$, there is a natural number $r$ such that the upper-tail
indicator of the normalized Rademacher sum on the $(r+1)$-dimensional sign cube, at the
Gaussian upper-tail quantile corresponding to mass $1/(q+2)$, has mean within half that
mass of $1/(q+2)$. Its total Fourier weight at level one is also within half the square
of the Gaussian density at that quantile of that squared density.
\end{theorem}

\begin{definition}[Selected approximation dimension]
\lean{FABL.fknOptimalityDimensionIndex}
\label{FABL.fknOptimalityDimensionIndex}
\leanok
\uses{FABL.exists_fknOptimalityApproximation}
For each natural number $q$, define an approximation dimension by selecting a natural
number $r$ for which the approximation property holds for $q$ and $r$.
\end{definition}

\begin{theorem}[Approximation property of the chosen dimension]
\lean{FABL.fknOptimalityDimensionIndex_spec}
\label{FABL.fknOptimalityDimensionIndex_spec}
\leanok
\uses{FABL.IsFKNOptimalityApproximation, FABL.exists_fknOptimalityApproximation, FABL.fknOptimalityDimensionIndex}
\difficulty{1}
For each natural number $q$, put $\alpha=1/(q+2)$ and let $t$ be the standard Gaussian
upper-tail quantile of mass $\alpha$. There is a dimension $N$ such that the upper-tail
indicator of the normalized sum of $N$ uniform signs has mean within $\alpha/2$ of
$\alpha$, and its level-one Fourier weight is within half the square of the standard
Gaussian density at $t$ of that square; the dimension index chosen for $q$, plus one,
has this property.
\end{theorem}

\begin{definition}[Tail dimension]
\lean{FABL.fknOptimalityTailDimension}
\label{FABL.fknOptimalityTailDimension}
\leanok
\uses{FABL.fknOptimalityDimensionIndex}
For each natural number $q$, define the tail dimension to be one greater than the
selected approximation dimension for $q$.
\end{definition}

\begin{definition}[Small-set indicator]
\lean{FABL.fknOptimalitySmallSet}
\label{FABL.fknOptimalitySmallSet}
\leanok
\uses{FABL.fknOptimalityTailDimension, FABL.fknOptimalityThreshold, FABL.hammingUpperTailIndicator}
For each natural number $q$, define a function on the sign cube of dimension equal to
the tail dimension for $q$: it is the Hamming upper-tail indicator at the threshold for
$q$.
\end{definition}

\begin{definition}[Small-set mass]
\lean{FABL.fknOptimalitySmallSetMass}
\label{FABL.fknOptimalitySmallSetMass}
\leanok
\uses{FABL.fknOptimalitySmallSet, FABL.fknOptimalityTailDimension}
For each natural number $q$, define the small-set mass as the expectation of its
indicator function under the uniform distribution on the sign cube of dimension equal to
the tail dimension for $q$.
\end{definition}

\begin{theorem}[Small-Set Mass Approximation]
\lean{FABL.fknOptimalitySmallSetMass_approx}
\label{FABL.fknOptimalitySmallSetMass_approx}
\leanok
\uses{FABL.fknOptimalityDimensionIndex, FABL.fknOptimalityDimensionIndex_spec, FABL.fknOptimalityGaussianMass, FABL.fknOptimalitySmallSetMass, FABL.fknOptimalityThreshold, FABL.hammingUpperTailIndicator}
\difficulty{3}
For a natural number $q$, set $\alpha=1/(q+2)$ and let $t$ be the standard Gaussian
upper-tail quantile of mass $\alpha$. In the Boolean cube of dimension $d$ selected by
the approximation property for $q$, let $A$ consist of the sign vectors whose normalized
coordinate sum exceeds $t$, and let $\mu(A)$ be its uniform probability. Then
$|\mu(A)-\alpha|<\alpha/2$.
\end{theorem}

\begin{theorem}[Approximation of Level-One Fourier Weight]
\lean{FABL.fknOptimalitySmallSetWeight_approx}
\label{FABL.fknOptimalitySmallSetWeight_approx}
\leanok
\uses{FABL.IsFKNOptimalityApproximation, FABL.fknOptimalityDimensionIndex_spec, FABL.fknOptimalityGaussianBoundary, FABL.fknOptimalitySmallSet, FABL.fknOptimalityTailDimension, FABL.fourierWeightAtLevel}
\difficulty{3}
For a natural number $q$, set $\alpha=1/(q+2)$, let $t$ be the standard Gaussian
upper-tail quantile of mass $\alpha$, and write $b$ for the Gaussian density at $t$. In
the Boolean cube of dimension $d$ selected by the approximation property for $q$, let
$A$ consist of the sign vectors whose normalized coordinate sum exceeds $t$. The sum of
the squared level-one Fourier coefficients of the indicator of $A$ differs from $b^2$ by
less than $b^2/2$.
\end{theorem}

\begin{theorem}[Positivity of Small-Set Mass]
\lean{FABL.fknOptimalitySmallSetMass_pos}
\label{FABL.fknOptimalitySmallSetMass_pos}
\leanok
\uses{FABL.fknOptimalityGaussianMass, FABL.fknOptimalitySmallSetMass, FABL.fknOptimalitySmallSetMass_approx}
\difficulty{3}
For a natural number $q$, set $\alpha=1/(q+2)$ and let $t$ be the standard Gaussian
upper-tail quantile of mass $\alpha$. In the Boolean cube of dimension $d$ selected by
the approximation property for $q$, let $A$ consist of the sign vectors whose normalized
coordinate sum exceeds $t$, and let $\mu(A)$ be its uniform probability. Then
$\mu(A)>0$.
\end{theorem}

\begin{theorem}[Gaussian Mass and Small-Set Mass]
\lean{FABL.fknOptimalityGaussianMass_lt_two_mul_smallSetMass}
\label{FABL.fknOptimalityGaussianMass_lt_two_mul_smallSetMass}
\leanok
\uses{FABL.fknOptimalityGaussianMass, FABL.fknOptimalitySmallSetMass, FABL.fknOptimalitySmallSetMass_approx}
\difficulty{3}
For a natural number $q$, set $\alpha=1/(q+2)$ and let $t$ be the standard Gaussian
upper-tail quantile of mass $\alpha$. In the Boolean cube of dimension $d$ selected by
the approximation property for $q$, let $A$ consist of the sign vectors whose normalized
coordinate sum exceeds $t$, and let $\mu(A)$ be its uniform probability. Then
$\alpha<2\mu(A)$.
\end{theorem}

\begin{theorem}[Small-set mass bound]
\lean{FABL.two_mul_fknOptimalitySmallSetMass_lt_three_mul_gaussianMass}
\label{FABL.two_mul_fknOptimalitySmallSetMass_lt_three_mul_gaussianMass}
\leanok
\uses{FABL.fknOptimalityGaussianMass, FABL.fknOptimalitySmallSetMass, FABL.fknOptimalitySmallSetMass_approx}
\difficulty{3}
For each natural number $q$, let $G_q=1/(q+2)$, and let $M_q$ be the uniform probability
mass of the associated small set in its Boolean cube. Then $2M_q<3G_q$.
\end{theorem}

\begin{theorem}[Lower Bound for Level-One Fourier Weight]
\lean{FABL.fknOptimalitySmallSetWeight_gt_half_boundary_sq}
\label{FABL.fknOptimalitySmallSetWeight_gt_half_boundary_sq}
\leanok
\uses{FABL.fknOptimalityGaussianBoundary, FABL.fknOptimalitySmallSet, FABL.fknOptimalitySmallSetWeight_approx, FABL.fourierWeightAtLevel}
\difficulty{3}
For a natural number $q$, set $\alpha=1/(q+2)$, let $t$ be the standard Gaussian
upper-tail quantile of mass $\alpha$, and write $b$ for the Gaussian density at $t$. In
the Boolean cube of dimension $d$ selected by the approximation property for $q$, let
$A$ consist of the sign vectors whose normalized coordinate sum exceeds $t$. The sum of
the squared level-one Fourier coefficients of the indicator of $A$ is greater than
$b^2/2$.
\end{theorem}

\begin{definition}[FKN optimality delta]
\lean{FABL.fknOptimalityDelta}
\label{FABL.fknOptimalityDelta}
\leanok
\uses{FABL.fknOptimalitySmallSetMass}
For each natural number $q$, define the FKN optimality delta to be twice the small-set
mass associated with $q$.
\end{definition}

\begin{theorem}[Positivity of the optimality parameter]
\lean{FABL.fknOptimalityDelta_pos}
\label{FABL.fknOptimalityDelta_pos}
\leanok
\uses{FABL.fknOptimalityDelta, FABL.fknOptimalitySmallSetMass_pos}
\difficulty{2}
For each natural number $q$, choose a dimension $N=r+1$ such that the upper-tail
indicator of the normalized sum of $N$ uniform signs has mass within half of $1/(q+2)$
of that value, and level-one Fourier weight within half the square of the corresponding
standard Gaussian density of that square. The parameter equal to twice the mass of this
upper-tail indicator is strictly positive.
\end{theorem}

\begin{theorem}[Gaussian mass below twice the set mass]
\lean{FABL.fknOptimalityGaussianMass_lt_delta}
\label{FABL.fknOptimalityGaussianMass_lt_delta}
\leanok
\uses{FABL.fknOptimalityDelta, FABL.fknOptimalityGaussianMass, FABL.fknOptimalityGaussianMass_lt_two_mul_smallSetMass}
\difficulty{1}
For every $q\in\mathbb{N}$, the Gaussian mass $1/(q+2)$ is strictly less than twice the
uniform measure of the associated small set.
\end{theorem}

\begin{theorem}[Upper bound for the optimality parameter]
\lean{FABL.fknOptimalityDelta_lt_three_mul_gaussianMass}
\label{FABL.fknOptimalityDelta_lt_three_mul_gaussianMass}
\leanok
\uses{FABL.fknOptimalityDelta, FABL.fknOptimalityGaussianMass, FABL.two_mul_fknOptimalitySmallSetMass_lt_three_mul_gaussianMass}
\difficulty{1}
For each natural number $q$, set $\alpha=1/(q+2)$ and let $t$ be the standard Gaussian
upper-tail quantile with upper-tail mass $\alpha$. Choose a dimension $N=r+1$ for which
the upper-tail indicator of the normalized sum of $N$ independent uniform signs has mass
within $\alpha/2$ of $\alpha$ and level-one Fourier weight within half the square of the
Gaussian density at $t$ of that square. If $\Delta$ is twice the mass of this upper-tail
indicator in dimension $N$, then $\Delta<3\alpha$.
\end{theorem}

\begin{theorem}[Vanishing Gaussian tail mass]
\lean{FABL.tendsto_fknOptimalityDelta}
\label{FABL.tendsto_fknOptimalityDelta}
\leanok
\uses{FABL.fknOptimalityDelta, FABL.fknOptimalityDelta_lt_three_mul_gaussianMass, FABL.fknOptimalityDelta_pos, FABL.tendsto_fknOptimalityGaussianMass_zero}
\difficulty{3}
For each natural number $q$, let $A_q$ be the proportion of points in the Boolean cube
of dimension $d_q$ whose normalized Rademacher sum exceeds the Gaussian upper-tail
quantile corresponding to mass $1/(q+2)$, where $d_q$ is the dimension chosen so that
this proportion approximates that mass and its level-one Fourier weight approximates the
square of the corresponding Gaussian boundary measure. Then $2A_q$ tends to $0$ as $q$
tends to infinity.
\end{theorem}

\begin{definition}[FKN optimality function]
\lean{FABL.fknOptimalityFunction}
\label{FABL.fknOptimalityFunction}
\leanok
\uses{FABL.BooleanFunction, FABL.fknOptimalitySmallSet, FABL.fknOptimalityTailDimension}
For each natural number $q$, consider inputs $z$ with signs in $\\{−1,1\\}$ and
dimension one greater than the FKN optimality tail dimension for $q$. The output is $1$
when the tail of $z$ lies in the small set associated with $q$; otherwise, the output is
the first coordinate of $z$.
\end{definition}

\begin{theorem}[Value on the positive first coordinate]
\lean{FABL.fknOptimalityFunction_toReal_cons_one}
\label{FABL.fknOptimalityFunction_toReal_cons_one}
\leanok
\uses{FABL.BooleanFunction.toReal, FABL.fknOptimalityFunction, FABL.fknOptimalityTailDimension}
\difficulty{1}
For every $q\in\mathbb{N}$ and every sign vector $x$ of the associated tail dimension,
the real-valued version of the constructed Boolean function, evaluated at the vector
whose first coordinate is $1$ and whose remaining coordinates are $x$, equals $1$.
\end{theorem}

\begin{theorem}[Value on the negative first coordinate]
\lean{FABL.fknOptimalityFunction_toReal_cons_neg_one}
\label{FABL.fknOptimalityFunction_toReal_cons_neg_one}
\leanok
\uses{FABL.BooleanFunction.toReal, FABL.fknOptimalityFunction, FABL.fknOptimalitySmallSet, FABL.fknOptimalityTailDimension, FABL.fknOptimalityThreshold, FABL.hammingUpperTailIndicator, FABL.normalizedRademacherSum}
\difficulty{4}
For every $q\in\mathbb{N}$ and every sign vector $x$ in the cube of dimension equal to
the tail dimension associated with $q$, the real-valued version of the constructed
Boolean function, evaluated at the vector whose first coordinate is $-1$ and whose
remaining coordinates are $x$, equals $2$ times the indicator of the associated small
set at $x$ minus $1$.
\end{theorem}

\begin{theorem}[Positive slice of the optimality function]
\lean{FABL.firstCoordinateSlice_fknOptimalityFunction_one}
\label{FABL.firstCoordinateSlice_fknOptimalityFunction_one}
\leanok
\uses{FABL.BooleanFunction.toReal, FABL.firstCoordinateSlice, FABL.fknOptimalityFunction, FABL.fknOptimalityFunction_toReal_cons_one}
\difficulty{3}
For every $q\in\mathbb N$, fixing the first coordinate of the real-valued sign embedding
of the optimality Boolean function to $1$ gives the constant function with value $1$ on
the tail coordinates.
\end{theorem}

\begin{theorem}[Negative slice of the optimality function]
\lean{FABL.firstCoordinateSlice_fknOptimalityFunction_neg_one}
\label{FABL.firstCoordinateSlice_fknOptimalityFunction_neg_one}
\leanok
\uses{FABL.BooleanFunction.toReal, FABL.firstCoordinateSlice, FABL.fknOptimalityFunction, FABL.fknOptimalityFunction_toReal_cons_neg_one, FABL.fknOptimalitySmallSet}
\difficulty{3}
For every $q\in\mathbb N$, fixing the first coordinate of the real-valued sign embedding
of the optimality Boolean function to $-1$ gives the function $x\mapsto 2A_q(x)-1$,
where $A_q$ is the small-set indicator on the tail coordinates.
\end{theorem}

\begin{theorem}[Empty-frequency coefficient of a centered indicator]
\lean{FABL.fourierCoeff_fknOptimalityLowerSlice_empty}
\label{FABL.fourierCoeff_fknOptimalityLowerSlice_empty}
\leanok
\uses{FABL.fknOptimalitySmallSet, FABL.fknOptimalitySmallSetMass, FABL.fourierCoeff, FABL.monomial}
\difficulty{3}
For each natural number $q$, let $A$ be the Gaussian-threshold upper-tail set in the
chosen tail dimension, and let $m$ be its uniform mass on that sign cube. The Fourier
coefficient at the empty frequency of the function $x\mapsto 2\mathbf{1}_A(x)-1$ is
$2m-1$.
\end{theorem}

\begin{theorem}[Singleton coefficient of a centered indicator]
\lean{FABL.fourierCoeff_fknOptimalityLowerSlice_singleton}
\label{FABL.fourierCoeff_fknOptimalityLowerSlice_singleton}
\leanok
\uses{FABL.expect_monomial, FABL.fknOptimalitySmallSet, FABL.fknOptimalityTailDimension, FABL.fourierCoeff, FABL.monomial}
\difficulty{4}
For each natural number $q$, let $A$ be the Gaussian-threshold upper-tail set in the
chosen tail dimension. For every coordinate $i$ in that dimension, the Fourier
coefficient at $i$ of $x\mapsto 2\mathbf{1}_A(x)-1$ equals twice the Fourier
coefficient at $i$ of the indicator of $A$.
\end{theorem}

\begin{theorem}[First-coordinate Fourier coefficient]
\lean{FABL.fourierCoeff_fknOptimalityFunction_zero}
\label{FABL.fourierCoeff_fknOptimalityFunction_zero}
\leanok
\uses{FABL.BooleanFunction.toReal, FABL.firstCoordinateSlice_fknOptimalityFunction_neg_one, FABL.firstCoordinateSlice_fknOptimalityFunction_one, FABL.fknOptimalityDelta, FABL.fknOptimalityFunction, FABL.fknOptimalityTailDimension, FABL.fourierCoeff, FABL.fourierCoeff_fknOptimalityLowerSlice_empty, FABL.fourierCoeff_insert_zero_tailFrequency, FABL.monomial, FABL.tailFrequency}
\difficulty{4}
For each natural number $q$, form the Boolean function on the sign cube of dimension one
greater than the chosen tail dimension: it is $1$ when the tail coordinates lie in the
associated Gaussian-threshold upper-tail set, and otherwise equals the first coordinate.
Its Fourier coefficient at the singleton frequency consisting of the first coordinate is
$1$ minus the mass of that set.
\end{theorem}

\begin{theorem}[Shifted Fourier coefficient identity]
\lean{FABL.fourierCoeff_fknOptimalityFunction_succ}
\label{FABL.fourierCoeff_fknOptimalityFunction_succ}
\leanok
\uses{FABL.BooleanFunction.toReal, FABL.expect_monomial, FABL.firstCoordinateSlice_fknOptimalityFunction_neg_one, FABL.firstCoordinateSlice_fknOptimalityFunction_one, FABL.fknOptimalityFunction, FABL.fknOptimalitySmallSet, FABL.fknOptimalityTailDimension, FABL.fourierCoeff, FABL.fourierCoeff_fknOptimalityLowerSlice_singleton, FABL.fourierCoeff_tailFrequency, FABL.tailFrequency}
\difficulty{4}
For each natural number $q$, let $d$ be the dimension chosen for the corresponding
Gaussian-threshold upper-tail set, and let $F$ be the Boolean function on the
$(d+1)$-dimensional sign cube that equals $1$ when the tail coordinates lie in that set
and otherwise equals the first coordinate. For every coordinate $i<d$, the Fourier
coefficient of the real-valued function associated to $F$ at the singleton frequency
containing coordinate $i+1$ equals the Fourier coefficient of the set's indicator at the
singleton frequency containing coordinate $i$.
\end{theorem}

\begin{theorem}[Fourier Weight of the Optimality Function]
\lean{FABL.fourierWeightAtLevel_one_fknOptimalityFunction}
\label{FABL.fourierWeightAtLevel_one_fknOptimalityFunction}
\leanok
\uses{FABL.BooleanFunction.toReal, FABL.fknOptimalityDelta, FABL.fknOptimalityFunction, FABL.fknOptimalitySmallSet, FABL.fknOptimalitySmallSetMass, FABL.fourierCoeff_fknOptimalityFunction_succ, FABL.fourierCoeff_fknOptimalityFunction_zero, FABL.fourierWeightAtLevel, FABL.fourierWeightAtLevel_one_eq_sum_singleton}
\difficulty{3}
For each natural number $q$, set $p_q=1/(q+2)$ and let $t_q$ be the standard Gaussian
upper-tail quantile with tail probability $p_q$. Choose a dimension $d_q$ for which the
indicator $A_q$ of the event that the normalized sum of the $d_q$ Boolean coordinates
exceeds $t_q$ has mean within $p_q/2$ of $p_q$ and level-one Fourier weight within half
the square of the Gaussian density at $t_q$ of that squared density. Let $F_q$ on the
$(d_q+1)$-dimensional cube equal the first coordinate when the remaining coordinates are
outside this event, and equal $1$ when they are inside it. Then $F_q$ has level-one
Fourier weight $(1-\mathbb{E}A_q)^2+W_1(A_q)$.
\end{theorem}

\begin{theorem}[Fourier Weight in Terms of the Set Mass]
\lean{FABL.fourierWeightAtLevel_one_fknOptimalityFunction_eq_delta}
\label{FABL.fourierWeightAtLevel_one_fknOptimalityFunction_eq_delta}
\leanok
\uses{FABL.BooleanFunction.toReal, FABL.fknOptimalityDelta, FABL.fknOptimalityFunction, FABL.fknOptimalitySmallSet, FABL.fourierWeightAtLevel, FABL.fourierWeightAtLevel_one_fknOptimalityFunction}
\difficulty{3}
For each natural number $q$, let $A_q$ be the upper-tail indicator on the selected
Boolean cube used to define the optimality function, and write
$\delta_q=2\mathbb{E}A_q$. Form $F_q$ on the cube with one additional coordinate by
setting it equal to $1$ when the remaining coordinates lie in the set and equal to the
first coordinate otherwise. Then $W_1(F_q)=1-\delta_q+\delta_q^2/4+W_1(A_q)$.
\end{theorem}

\begin{definition}[Gaussian FKN optimality quantity]
\lean{FABL.fknOptimalityGaussianMain}
\statementsource{Analysis-of-Boolean-Functions-by-Ryan-ODonnell}{
  pdf-ffbe9c7b3306-p061-b001,
  pdf-ffbe9c7b3306-p061-b002
}
\label{FABL.fknOptimalityGaussianMain}
\leanok
\uses{FABL.fknOptimalityGaussianMass}
For each natural number $q$, let $m_q$ denote the specified Gaussian mass at $q$. Define
the quantity $m_q\sqrt{2\log(1/m_q)}$.
\end{definition}

\begin{theorem}[Reciprocal Gaussian mass bound]
\lean{FABL.one_lt_inv_fknOptimalityGaussianMass}
\label{FABL.one_lt_inv_fknOptimalityGaussianMass}
\leanok
\uses{FABL.fknOptimalityGaussianMass}
\difficulty{3}
For every natural number $q$, let $\gamma_q=1/(q+2)$. Then $1<1/\gamma_q$.
\end{theorem}

\begin{theorem}[Positive Gaussian main term]
\lean{FABL.fknOptimalityGaussianMain_pos}
\label{FABL.fknOptimalityGaussianMain_pos}
\leanok
\uses{FABL.fknOptimalityGaussianMain, FABL.fknOptimalityGaussianMass, FABL.one_lt_inv_fknOptimalityGaussianMass}
\difficulty{3}
For every $q\in\mathbb{N}$, the quantity $\frac{1}{q+2}\sqrt{2\log(q+2)}$ is strictly
positive.
\end{theorem}

\begin{theorem}[Asymptotic Gaussian boundary equivalence]
\lean{FABL.fknOptimalityGaussianBoundary_isEquivalent_main}
\label{FABL.fknOptimalityGaussianBoundary_isEquivalent_main}
\leanok
\uses{FABL.fknOptimalityGaussianBoundary, FABL.fknOptimalityGaussianMain, FABL.gaussianIsoperimetric_isEquivalent_atBot, FABL.tendsto_fknOptimalityGaussianMass_atBot}
\difficulty{1}
As $q\to\infty$ through the natural numbers, the Gaussian boundary associated with $q$
is asymptotically equivalent to $\frac{1}{q+2}\sqrt{2\log(q+2)}$; that is, their ratio
tends to $1$.
\end{theorem}

\begin{theorem}[Gaussian main term below boundary]
\lean{FABL.eventually_half_fknOptimalityGaussianMain_lt_boundary}
\label{FABL.eventually_half_fknOptimalityGaussianMain_lt_boundary}
\leanok
\uses{FABL.fknOptimalityGaussianBoundary, FABL.fknOptimalityGaussianBoundary_isEquivalent_main, FABL.fknOptimalityGaussianMain, FABL.fknOptimalityGaussianMain_pos}
\difficulty{4}
For all sufficiently large integers $q$, half of $m_q\sqrt{2\log(1/m_q)}$ is strictly
less than $b_q$, where $m_q=1/(q+2)$ is the prescribed Gaussian mass and $b_q$ is its
Gaussian isoperimetric boundary value.
\end{theorem}

\begin{theorem}[Square of the Gaussian main term]
\lean{FABL.fknOptimalityGaussianMain_sq}
\label{FABL.fknOptimalityGaussianMain_sq}
\leanok
\uses{FABL.fknOptimalityGaussianMain, FABL.fknOptimalityGaussianMass, FABL.one_lt_inv_fknOptimalityGaussianMass}
\difficulty{3}
For every $q\in\mathbb{N}$, the square of $\frac{1}{q+2}\sqrt{2\log(q+2)}$ equals
$2\bigl(\frac{1}{q+2}\bigr)^2\log(q+2)$.
\end{theorem}

\begin{theorem}[Gaussian boundary asymptotic bound]
\lean{FABL.eventually_gaussianMass_sq_log_div_two_lt_boundary_sq}
\label{FABL.eventually_gaussianMass_sq_log_div_two_lt_boundary_sq}
\leanok
\uses{FABL.eventually_half_fknOptimalityGaussianMain_lt_boundary, FABL.fknOptimalityGaussianBoundary, FABL.fknOptimalityGaussianBoundary_pos, FABL.fknOptimalityGaussianMain, FABL.fknOptimalityGaussianMain_pos, FABL.fknOptimalityGaussianMain_sq, FABL.fknOptimalityGaussianMass}
\difficulty{4}
For all sufficiently large integers $q$, $\frac{m_q^2}{2}\log(1/m_q)<b_q^2$, where
$m_q=1/(q+2)$ is the prescribed Gaussian mass and $b_q$ is its Gaussian isoperimetric
boundary value.
\end{theorem}

\begin{theorem}[Small-set Fourier weight lower bound]
\lean{FABL.eventually_gaussianMass_sq_log_div_four_lt_smallSetWeight}
\label{FABL.eventually_gaussianMass_sq_log_div_four_lt_smallSetWeight}
\leanok
\uses{FABL.eventually_gaussianMass_sq_log_div_two_lt_boundary_sq, FABL.fknOptimalityGaussianMass, FABL.fknOptimalitySmallSet, FABL.fknOptimalitySmallSetWeight_gt_half_boundary_sq, FABL.fourierWeightAtLevel}
\difficulty{3}
For all sufficiently large integers $q$, the level-one Fourier weight of the small-set
indicator exceeds $\frac{m_q^2}{4}\log(1/m_q)$, where $m_q=1/(q+2)$ is the prescribed
Gaussian mass.
\end{theorem}

\begin{theorem}[Logarithmic bound for the optimality parameter]
\lean{FABL.fknOptimalityDelta_sq_log_le}
\label{FABL.fknOptimalityDelta_sq_log_le}
\leanok
\uses{FABL.fknOptimalityDelta, FABL.fknOptimalityDelta_lt_three_mul_gaussianMass, FABL.fknOptimalityDelta_pos, FABL.fknOptimalityGaussianMass, FABL.fknOptimalityGaussianMass_lt_delta}
\difficulty{5}
For each natural number $q$, let $\alpha=1/(q+2)$, and let $\Delta$ be twice the mass of
the upper-tail indicator of the normalized sum of uniform signs in a dimension chosen so
that its mass approximates $\alpha$ and its level-one Fourier weight approximates the
square of the corresponding standard Gaussian density, as specified by the two error
bounds. If $\Delta<1$, then $\Delta^2\log(1/\Delta)\leq 9\alpha^2\log(1/\alpha)$.
\end{theorem}

\begin{theorem}[Eventual smallness of the doubled set mass]
\lean{FABL.eventually_fknOptimalityDelta_lt_one}
\label{FABL.eventually_fknOptimalityDelta_lt_one}
\leanok
\uses{FABL.fknOptimalityDelta, FABL.tendsto_fknOptimalityDelta}
\difficulty{1}
For all sufficiently large integers $q$, the doubled uniform mass of the threshold set
in its selected dimension is less than $1$.
\end{theorem}

\begin{theorem}[Eventual Fourier weight gain]
\lean{FABL.eventually_fknOptimalityGain}
\label{FABL.eventually_fknOptimalityGain}
\leanok
\uses{FABL.eventually_fknOptimalityDelta_lt_one, FABL.eventually_gaussianMass_sq_log_div_four_lt_smallSetWeight, FABL.fknOptimalityDelta, FABL.fknOptimalityDelta_sq_log_le, FABL.fknOptimalitySmallSet, FABL.fourierWeightAtLevel}
\difficulty{4}
For all sufficiently large integers $q$, the level-one Fourier weight of the small-set
indicator is at least $\frac{1}{36}\Delta_q^2\log(1/\Delta_q)$, where $\Delta_q$ is
twice its uniform mass.
\end{theorem}

\begin{theorem}[Eventual Fourier weight of the modified function]
\lean{FABL.eventually_fourierWeightAtLevel_one_fknOptimalityFunction_lower}
\label{FABL.eventually_fourierWeightAtLevel_one_fknOptimalityFunction_lower}
\leanok
\uses{FABL.BooleanFunction.toReal, FABL.eventually_fknOptimalityGain, FABL.fknOptimalityDelta, FABL.fknOptimalityFunction, FABL.fourierWeightAtLevel, FABL.fourierWeightAtLevel_one_fknOptimalityFunction_eq_delta}
\difficulty{3}
For all sufficiently large integers $q$, the level-one Fourier weight of the real-valued
version of the Boolean function obtained by using the small-set indicator on the tail
coordinates and otherwise returning the first coordinate is at least
$1-\Delta_q+\frac{1}{36}\Delta_q^2\log(1/\Delta_q)$, where $\Delta_q$ is twice the
uniform mass of the small set.
\end{theorem}

\begin{theorem}[Near-extremal first-level Fourier weight]
\lean{FABL.exercise5_36}
\label{FABL.exercise5_36}
\leanok
\uses{FABL.BooleanFunction, FABL.BooleanFunction.toReal, FABL.eventually_fourierWeightAtLevel_one_fknOptimalityFunction_lower, FABL.fknOptimalityDelta, FABL.fknOptimalityDelta_pos, FABL.fknOptimalityFunction, FABL.fknOptimalityTailDimension, FABL.fourierCoeff, FABL.fourierCoeff_fknOptimalityFunction_zero, FABL.fourierWeightAtLevel, FABL.tendsto_fknOptimalityDelta}
\difficulty{3}
There exist a sequence of natural numbers $d_q$, Boolean functions
$f_q:\{-1,1\}^{d_q+1}\to\{-1,1\}$, and positive real numbers $\delta_q$ tending to $0$
as $q\to\infty$, such that the Fourier coefficient of $f_q$ at the singleton coordinate
set $\{0\}$ is $1-\delta_q/2$ for every $q$. Moreover, there is a constant $c>0$ such
that, for all sufficiently large $q$,
\[
1-\delta_q+c\delta_q^2\log(1/\delta_q)\le W_1(f_q),
\]
where $W_1(f_q)$ is the sum of the squares of the Fourier coefficients over all
singleton coordinate sets.
\end{theorem}
